\section{Founder Adaptation：从确定性组织到低概率系统}

课程随后转向创始人角色。McKinnon 离开 Salesforce 时，失去团队、品牌、行政支持和高概率晋升路径；创业系统的输入变成有限 runway、模糊 buyer signal 和不断变化的假设。本节保留这段 teacher voice，因为它解释了基础设施公司早期如何在没有完整证据时仍建立可验证进展。

\subsection{Resource Reset 与 daily evidence}

创始人不能把“信念”当作拒绝数据。更稳健的循环是把大愿景拆成每天可更新的 evidence：谁愿意接受访谈，哪种 pain 会触发预算，prototype 是否接入真实 directory，first tenant 是否稳定，第一笔付款是否重复。每一步都降低一类不确定性，并为下一轮招人、融资或产品取舍提供依据。

\lecturefigure{06-founder-adaptation.png}{从成熟公司到创业团队意味着资源和确定性归零，再用 daily evidence 重建执行系统。}{本地字幕 06:00--12:20；概念重绘。}

\begin{knowledgebox}{读图：Belief 与 probability 可以同时存在}
Resource reset 说明旧组织优势不能自动带走；daily evidence 把愿景转成 customer/prototype 信号；low odds 要求明确 runway 和 stopping rule；team belief 则需要方向、诚实与可见进展。讲者所说“即使不相信也要相信”不应理解为欺骗团队，而是领导者在承认概率不利时仍能组织行动。
\end{knowledgebox}

\teachervoice{McKinnon 坦率地说，离职的动机既有技术机会，也有“想自己当老板”的个人愿望；真正困难的是在金融危机、新生儿和家庭风险面前做共同决定。课堂把 founder risk 从个人英雄叙事拉回家庭、资源和机会成本。}

\begin{warningbox}{不要把 survivorship bias 写成创业公式}
Okta 后来的成功不能证明“离开大公司”“坚持足够久”或“中年经验”必然带来结果。可迁移原则是选择与自身能力匹配的 market、尽快获得 buyer evidence、限制不可逆风险，并在新证据出现时更新计划；结果仍受到时机、竞争、资本和运气影响。
\end{warningbox}

\subsection{本章小结}

Founder adaptation 是从 inherited resources 转向 evidence loop。信念负责持续行动，概率与指标负责纠偏；两者缺一都会让组织失去现实感。

